\scanpage{31}
\deleted{\unsure{VIII|IX|VII}} \uline{IX}

1. Suite des faisceaux de torsion

2. Faisceaux loc. constants.

3. Faisceaux constructibles

\uline{\unsure{Préciser|Précise} nos conventions}. Qu'est-ce que veut dire loc. ct. à fibres finies sur $X^{\mathrm{cons}}$. Le reste du n° 2 donne les propriétés de commutation de la notion

\textcircled{a} \uline{Conditions équivalentes}\,: (i) L'image inverse sur $X^{\unsure{\mathrm{ens}|\mathrm{cons}|\mathrm{red}}}$ est loc. ct à fibres finies

(ii) $\exists\ f:X'\unsure{-|\to}X$ loc. de présentation finie surjectif, avec $f^*(F)$ loc. ct à fibres finies

(iii) \deleted{Si} Si $X$ qu. cp et qu. sép) $\exists$ partition de $X$ en $X_i$ loc. fermés constructibles, en nb fini, t.q. $F|X_i$ loc. ct. à fibres finies.

(iv) \inserted{(Si $X$ qu. cp et qu. sép)} $F$ est le conoyau d'une double flèche
\[
\deleted{\unsure{F_1}} F_{\unsure{1|2}}\rightrightarrows F_0,
\]
où $F_0$ et $F_{1}$ sont des faisceaux repr. par des schémas étales de prés. finie sur $X$

(v) \inserted{($X$ qu. cp et qu. sép)} $F$ est \deleted{du type} \inserted{coker d'} \unsure{constructiblement|constructibles} de type fini, et tt \unsure{hom} d'un faisceau \deleted{\unsure{repr}} $G$ dans $F$, avec $G$ repr. par pr. étale de prés finie sur $X$, \deleted{\unsure{le}} la \unsure{rel.} équiv. est constr.

(vi) \inserted{($X$ qu. cp et qu. sép\,; Gr. abéliens \deleted{\unsure{$F$}} et \unsure{cat})} $F$ admet une filtration par des $i_!(G)$, $i:Y-X$ imm. \inserted{constr.} \deleted{qu. cp \unsure{et}}, $G$ \deleted{\unsure{repr est}} sur $Y$

(vii) ($X$ qu. cp et qu. sép.\,, Gr abéliens, $A$ et \unsure{autres})\,:\enspace $F$ admet une \deleted{\unsure{prés}} résolution \inserted{à deux crans} par des \unsure{$p_* (\mathbb Z_{X'})$|$p_*(G_{X'})$}, $X'$ fini de prés finie sur $X$, $G$ ct. \unsure{constr.}

% slanted
\margin{Pour \unsure{dévisser en|l'énoncé} \unsure{catégorie|algèbre} les faisceaux de torsion}

\inserted{Si $A=\mathbb Z$, faisceaux de l-torsion, $X$ qu. cp et qu. sép)} (viii) \deleted{\unsure{Si $F$ est constructible et $X$ localement} \illegible{3} \unsure{volumes} \illegible{2}} $F$ se dévisse par \deleted{\unsure{faisceaux}} \inserted{suite de cinq} \unsure{termes} en faisceaux \inserted{finis} dont des quotients facteurs directs d'un faisceau lié, \unsure{admettant} une suite de composition par des faisceaux de la forme $p_*(i_!(\mathbb Z/\ell\mathbb Z)_U)$, où $p:X'-X$ mo. fini surj., fini \inserted{et de} \deleted{\unsure{prés}} prés., $U$ ouvert \unsure{dense} de $X'$, $i:U-X'$ l'inclusion

On peut \unsure{couper|songer} avec des complexes de $\mathcal A$-Modules et la \unsure{$K$|$A$}-constructibilité

Au vu de ces rédactions, il sera bon de convenir pour la \unsure{$0$|$a$|$\mathbb Q$}- et la $1$-constructibilité.

\textcircled{b} \uline{Déf} Constructibles de type fini, resp. constr. de prés finie = constructibles

\deleted{et \unsure{constr. égaux}}

\uline{Cor} \textcircled{c} Stabilité par $\underset{\longrightarrow}{\lim}$ et $\underset{\longleftarrow}{\lim}$ finies (dans le cas abélien, il faut \deleted{et \unsure{eff}} que $A$ \unsure{noeth. resp.|soit un corps.} $\mathcal A$ soit \unsure{ann. cohérent}), pr \uline{facteurs} \uline{directs}.

\uline{Cor} \textcircled{d} Caractère épais (\unsure{Grs comm.|Grs non-comm.|cas noethérien})

\uline{Cor} \textcircled{e} Stabilité par image inverse. par image directe par morph fini de prés. finie \deleted{par image directe par} \deleted{\unsure{implication finie de prés. finie}}

% archive sheet number: 30
